\documentclass{article}
\usepackage[margin=1in]{geometry}
\usepackage{array}
\usepackage{longtable}
\usepackage{booktabs}
\usepackage{float}
\begin{document}
\footnotesize
\setlength{\tabcolsep}{4pt}
\renewcommand{\arraystretch}{0.9}
\section*{6.3 (B)}

\begin{table}[H]
\centering
\begin{tabular}{>{\raggedright\arraybackslash}p{0.20\textwidth} >{\raggedleft\arraybackslash}r >{\raggedright\arraybackslash}p{0.45\textwidth} >{\raggedleft\arraybackslash}r >{\raggedleft\arraybackslash}r}
\toprule
variable & dtype & label & missing & \% missing \\
\midrule
hhsize & float32 & household size & 0 & 0.00 \\
povrat & float32 & income to needs ratio & 628 & 8.07 \\
bmi & float32 & own body mass index (wt for height measure) & 1456 & 18.70 \\
ageyrs & int8 & age in years & 0 & 0.00 \\
female & int8 & =1 if female & 0 & 0.00 \\
white & int8 & =1 if white & 0 & 0.00 \\
black & int8 & =1 if black & 0 & 0.00 \\
hisp & int8 & =1 if hispanic & 0 & 0.00 \\
other & int8 & =1 if other race & 0 & 0.00 \\
tvyest & float64 & \# hours of TV watched yesterday & 1435 & 18.43 \\
dadbmi & float32 & dad's body mass index & 1256 & 16.13 \\
mombmi & float32 & mom's body mass index & 617 & 7.93 \\
dadobese & float64 & =1 if dad obese & 1256 & 16.13 \\
dadovwt & float64 & =1 if dad overweight & 1256 & 16.13 \\
momobese & float64 & =1 if mom obese & 617 & 7.93 \\
momovwt & float64 & =1 if mom overweight & 617 & 7.93 \\
ovwtc & float64 & =1 if overweight & 1456 & 18.70 \\
obesec & float64 & =1 if obese & 1456 & 18.70 \\
pctcal & float32 & calories consumed/2000 & 530 & 6.81 \\
\bottomrule
\end{tabular}
\end{table}

\clearpage
\section*{6.3 (C)}

\begin{table}[H]
\centering
\begin{tabular}{>{\raggedright\arraybackslash}p{0.24\textwidth} >{\raggedleft\arraybackslash}r >{\raggedleft\arraybackslash}r >{\raggedleft\arraybackslash}r >{\raggedleft\arraybackslash}r >{\raggedleft\arraybackslash}r}
\toprule
variable & N & mean & sd & min & max \\
\midrule
obesec & 6329 & 0.1343 & 0.3410 & 0.0000 & 1.0000 \\
ageyrs & 7785 & 9.4145 & 3.3664 & 5.0000 & 16.0000 \\
female & 7785 & 0.5070 & 0.5000 & 0.0000 & 1.0000 \\
white & 7785 & 0.3188 & 0.4660 & 0.0000 & 1.0000 \\
black & 7785 & 0.3105 & 0.4627 & 0.0000 & 1.0000 \\
hisp & 7785 & 0.3166 & 0.4652 & 0.0000 & 1.0000 \\
tvyest & 6350 & 3.1411 & 1.8018 & 0.0000 & 6.0000 \\
povrat & 7157 & 1.7303 & 1.3162 & 0.0000 & 7.5170 \\
\bottomrule
\end{tabular}
\end{table}

\clearpage
\section*{6.3 (D)}

\renewcommand{\arraystretch}{1.1}
\setlength{\tabcolsep}{6pt}
\begin{table}[H]
\centering
\begin{tabular}{lrrrr}
\hline
Variable & Coef & Std Err & $t$ & $p$ \\
\hline
const & -1.560$^{\ast\ast\ast}$ & 0.131 & -11.917 & 0.000 \\
 & $(0.131)$ & & & \\
ageyrs & 0.008 & 0.006 & 1.267 & 0.205 \\
 & $(0.006)$ & & & \\
female & 0.030 & 0.042 & 0.706 & 0.480 \\
 & $(0.042)$ & & & \\
white & 0.024 & 0.116 & 0.204 & 0.838 \\
 & $(0.116)$ & & & \\
black & 0.204$^{\ast}$ & 0.112 & 1.828 & 0.068 \\
 & $(0.112)$ & & & \\
hisp & 0.261$^{\ast\ast}$ & 0.112 & 2.340 & 0.019 \\
 & $(0.112)$ & & & \\
tvyest & 0.057$^{\ast\ast\ast}$ & 0.012 & 4.789 & 0.000 \\
 & $(0.012)$ & & & \\
povrat & 0.002 & 0.018 & 0.108 & 0.914 \\
 & $(0.018)$ & & & \\
\hline
$N$ & 5770 & & & \\
Pseudo $R^2$ & 0.012 & & & \\
Log-likelihood & -2234.689 & & & \\
LR stat & 53.341 & & & \\
LR p-value & 0.000 & & & \\
AIC & 4485.379 & & & \\
BIC & 4538.662 & & & \\
\hline
\end{tabular}
\end{table}
\vspace{2pt}
\begin{flushleft}\footnotesize{* $p<0.1$, ** $p<0.05$, *** $p<0.01$}\end{flushleft}

\clearpage
\section*{6.3 (E)}

\noindent In the Probit model, the coefficient on \texttt{hisp} is positive and statistically significant (\(p=0.019\)), so Hispanic children have a higher latent propensity to be obese, holding other covariates fixed. The average marginal effect is about \(0.055\) (SE \(0.024\), \(p=0.019\)), implying roughly a 5.5 percentage-point higher probability of obesity on average.

\clearpage
\section*{6.3 (F)}

\noindent Using the fitted Probit model, the predicted probability of obesity for a 16-year-old female white child who watched 3 hours of TV yesterday with income-to-needs ratio 4.56 is \(\hat{p}=0.114\) (about 11.4\%).

\clearpage
\section*{6.3 (G)}

\noindent If the same child watched no TV yesterday, the predicted probability falls to \(\hat{p}=0.084\) (8.4\%). The drop is \(\Delta \hat{p}=0.0295\) (2.95 percentage points). Using a delta-method test for the difference gives SE \(0.0062\), \(z=4.78\), and \(p\approx1.8\times10^{-6}\), so the fall is statistically significant.

\end{document}
